\documentclass[sigconf,nonacm]{acmart}

\usepackage{booktabs}
\usepackage{listings}
\usepackage{xcolor}

\definecolor{rulegreen}{HTML}{1C7A4E}
\definecolor{ruleink}{HTML}{14181A}
\definecolor{rulefaint}{HTML}{7B868A}

\lstset{
  basicstyle=\ttfamily\footnotesize\color{ruleink},
  commentstyle=\color{rulefaint}\itshape,
  keywordstyle=\color{rulegreen}\bfseries,
  morekeywords={not},
  columns=fullflexible, keepspaces=true, breaklines=true,
  frame=leftline, framesep=6pt, rulecolor=\color{rulefaint},
  xleftmargin=10pt,
  literate={<-}{{$\,\leftarrow\,$}}{3}
}

\begin{document}

\title{An Executable Decision-Rights Matrix and an Expert-Approved Logic
Database}
\subtitle{Datalog with stratified negation for deterministic, explainable and
attributable decisions in human-centred agentic systems}

\author{confused.de}
\affiliation{%
  \institution{ACM Europe School on Responsible AI}
  \country{}
}

\begin{abstract}
Human-centred agentic AI is usually specified as a person reviewing each
consequential action. At volume that is unaffordable, so it is skipped or
performed at a rate that makes it ceremonial; and a reviewer shown a
recommendation without its reasoning can do little besides agree. The control
degrades into the appearance of control.

We report a working system built on a different allocation. Our primary
contribution is a \emph{decision-rights matrix made executable}, together with a
human-in-the-loop mechanism whose single question is not ``how critical is
this?'' but ``has a person already ruled on this logic?''. The matrix assigns one
of four control tiers from properties of the decision --- harm, reversibility,
rights, evidence --- and an expert reviews a \emph{prepared answer} they may
edit, not a bare question. Their ruling is then captured as logic, so the same
judgement is not demanded of them again.

That is only honest if the captured logic behaves. Our second contribution is
the substrate that makes it so: a logic database of function-free Horn clauses
with three-valued facts, in which the same case yields the same answer and the
same derivation, every conclusion carries the proof that produced it, and every
rule carries the name of whoever approved it.

Over a twelve-case load from an empty database, 18--30 per cent of decisions were
made deterministically from approved logic. We report what broke in building it,
and what the system still cannot do.
\end{abstract}

\maketitle

\section{Problem: per-decision oversight does not scale}

Oversight of agentic systems is conventionally specified per action. A person
reviews each consequential output before it takes effect, and the arrangement is
called meaningful human control.

It fails in two directions at once. At volume, per-decision review is
unaffordable, so it is either skipped or performed at a rate that makes it
ceremonial. And a reviewer shown a recommendation without the reasoning that
produced it cannot do much besides agree: the check that was meant to catch
errors becomes a signature that launders them.

The usual response is to tune the threshold --- review the top few per cent,
sample, escalate on low confidence. That manages the cost without addressing the
cause. Our claim is that human control scales only if the unit of human
judgement is the \emph{rule} rather than the case, and that this requires two
things to be separated which are normally conflated: how much human involvement
a decision warrants, and what the decision actually is.

\section{Contribution 1: an executable decision-rights matrix}

\subsection{From written specification to executable classifier}

Our group's earlier prototype produced a decision-rights matrix: four tiers of
AI authority assigned from harm, reversibility, whether a right is engaged, and
how complete the evidence is. It is a restatement of the structure running
through the EU AI Act and the assurance literature, and like most such artefacts
it lived in a document.

We made it executable. The four tiers are evaluated top down:

\begin{lstlisting}
if (harm >= 9 and rights) or (harm >= 8 and reversibility <= 3):
    human_only          # AI may not decide
if harm >= 6 or rights or confidence < 70:
    human_approval      # a person approves before it takes effect
if harm <= 3 and reversibility >= 7:
    ai_with_oversight   # bounded autonomy, logged and reversible
ai_recommends           # analyse and propose; do not act
\end{lstlisting}

Nine lines, in a module that imports nothing --- not the graph framework, not
the model client, not the rest of the package. A module that cannot reach a
network cannot be talked into anything, and this is the one that decides how
much supervision a decision receives.

The classification also runs a second time, in JavaScript, so the console's
sliders respond without a round trip. Two copies of the rule allocating human
oversight is a hazard: they can drift with nothing failing, and the failure mode
is a page showing one tier whilst the server acts on another. A test extracts
the classifier \emph{from the shipped page}, runs it under Node, and compares all
1{,}000 input combinations against the Python. Mutating one threshold produced
twenty disagreements, each named.

A language model estimates the four inputs from the conversation and returns
nothing else. It is told explicitly that it does not decide what oversight
follows from its numbers. Estimating harm is a judgement about the world, which
a model may reasonably attempt; deciding what supervision follows is a policy,
which must be stable. The scales the model is given and the help text a person
reads behind each control are generated from one source, because two definitions
of ``harm'' make the two numbers look comparable when they are not.

\subsection{The oversight gate}

\begin{table}[t]
\caption{What the gate does. The tier sets what is at stake, not whether a
person is asked.}
\label{tab:gate}
\small
\begin{tabular}{@{}lll@{}}
\toprule
\textbf{Tier} & \textbf{Approved rule fires} & \textbf{Nothing fires} \\
\midrule
\texttt{human\_only}        & proceeds & \textbf{waits} \\
\texttt{human\_approval}    & proceeds & \textbf{waits} \\
\texttt{ai\_with\_oversight}& proceeds & proceeds, sampled \\
\texttt{ai\_recommends}     & proceeds & proceeds, sampled \\
\bottomrule
\end{tabular}
\end{table}

That approved logic proceeds even at \texttt{human\_only} is the design's most
contestable claim, and it is deliberate. A rule an expert approved \emph{is}
their decision: they saw the clause, they accepted that any case matching its
body receives its conclusion, and applying it executes a call they have already
made. Requiring them to confirm it again for every matching case is precisely
how oversight becomes ceremony.

Three things keep that defensible. Approval is of a \emph{clause}, whose body
states exactly which cases it claims. A proportion of decisions is audited
afterwards regardless, at a rate highest where a wrong decision hurts most ---
50 per cent for \texttt{human\_only} logic, 5 per cent for routine --- and never
zero, because a rule right a thousand times can begin to be wrong when the
population shifts beneath it. Sampling is deterministic on a hash of the case,
so an auditor can recompute why any case was or was not checked. And an approved
rule can be retired, which enumerates every decision it already made.

\subsection{Expert review of prepared decisions}

An inquiry that cannot proceed parks in the database rather than in the
framework's memory --- an in-memory checkpoint loses every pending approval on
restart, leaving the person who was waiting waiting for something that no longer
exists.

Crucially, the agent still \emph{prepares the answer}. The matrix says the AI
prepares a consequential action and a person rules on it; an expert handed a
bare question has nothing to rule on. So the draft is written, withheld from the
person waiting, and shown to the expert in an editable field. An expert who may
only accept or refuse is not reviewing, they are signing; the ability to change
the words is what makes the approval theirs. The answer finally sent is
attributed to them and says whether they revised it.

Their ruling is then put back as logic. A clause capturing what they decided is
drafted, checked that it fires on the case's own facts, and offered in the same
interaction:

\begin{lstlisting}
not_evict(c) <- in_arrears(c), disputes_debt(c),
                not formal_notice_given(c).
\end{lstlisting}

\noindent It is not approved automatically. They ruled on an \emph{answer}; the
clause is a generalisation of that answer and may be wider than they meant.
Approving it is one further click, made with the clause in front of them.

\section{Contribution 2: a Datalog logic database as formal substrate}

Capturing a person's judgement as logic is only honest if the logic behaves
predictably. That is what the second contribution provides.

The formalism is Datalog: function-free Horn clauses with stratified negation,
over facts carrying three truth values. Stratification is checked when the base
is loaded, so a clause set whose negations are mutually recursive is rejected
with the offending cycle named rather than evaluated to an arbitrary fixpoint.
Clause identity is canonical --- variables renumbered, body sorted --- and
safety (every head variable bound by a positive body literal) is checked at
construction, so an unsafe clause cannot enter the base. The properties below
are consequences of that choice rather than features added on top of it.

\textbf{Determinism.} Function-free Horn clauses evaluated by forward chaining
to a fixpoint. No function symbols, so the derivable set is finite and
termination is structural rather than a timeout. Each round matches a snapshot,
so a rule's result cannot depend on where in the round it ran, and insertion
order cannot change the answer.

\textbf{Explainability.} Forward rather than backward chaining, so the whole
derivation exists when evaluation finishes and the explanation is a transcript
of what happened rather than an account reconstructed afterwards. That
distinction is the difference between a proof and a story.

\textbf{Transparency.} Every rule carries its origin, its proposer, its
approver, a usage count and an audit tally. Every decision carries the facts it
was given, the rules that fired, and its tier. Rule history is append-only,
because the question an audit asks is not what the base says now but who changed
it and when.

\textbf{Evidence that stays honest.} Facts are three-valued, and the third value
carries the weight. Facts arrive from a model reading free text; if it did not
mention a formal notice, that does not mean none was served --- it means nobody
said. Under a closed world that silence becomes \texttt{false}, the negation
succeeds, and a rule fires against a person on evidence nobody gathered:

\begin{lstlisting}
will_get_wet(C) <- going_out(C), raining(C),
                   not has_umbrella(C).
\end{lstlisting}

\noindent With the umbrella recorded absent, the rule fires. With the umbrella
\emph{unmentioned}, it does not, and the head is recorded as blocked rather than
false, so anything reading \texttt{not will\_get\_wet} sees \textsc{unknown} too.
We shipped that defect twice in a predecessor system.

\textbf{What almost fired.} A clause needing three conditions with two satisfied
is not a miss but a question, and a precise one. The system reports it: \emph{is
\texttt{disputes\_debt(c)} the case?} Establishing one fact then settles the
matter by approved logic rather than by a model's judgement.

\textbf{Generalisation, with the consequences attached.} Two approved clauses
alike but for one condition each merge to their shared core; two differing only
in a value merge by variabilising it. Each candidate is replayed over every
stored case and presented with what it keeps, what it newly decides and what it
contradicts. Nothing is applied automatically: a generalisation fires on
strictly more cases than its sources, so it always risks deciding something
nobody intended.

\section{Failures encountered during development}

The system was built against a live model throughout. Five failures are worth
reporting.

\textbf{A greeting was classified.} Typing ``Hi'' produced \emph{human approval
required}. The estimator correctly scored confidence 0 and wrote ``no decision
has been described yet''; the ladder correctly read \texttt{confidence < 70} and
demanded a person. The architecture forced a classification out of a
conversation. The remedy was a separate intake phase that classifies nothing.

\textbf{A gate that overfitted its own prompt.} We first judged the intake
handover flaky. Measured, it was deterministically wrong: 0/6 on one ordinary
case and 6/6 on a near-identical one, the only difference being that the second
appeared in the prompt's example list. Four rounds of rewording produced
inconsistent results, which was the signal to change shape rather than words.
The gate now asks the \emph{narrow} question --- was this purely
conversational? --- and hands over unless the answer is yes, failing open
because being wrong that way costs a moment of needless assessment whilst being
wrong the other way means a decision is never assessed at all.

\textbf{The model was the limit, not the prompt.} Fact extraction returned
nothing on ordinary decisions. We spent four rounds on wording before measuring
the obvious alternative: on four cases, Claude Haiku 4.5 extracted facts from
\textbf{0}, and Sonnet from \textbf{4}. Extraction and rule drafting now take a
capability floor rather than the selected model.

\textbf{Conclusions were stored as facts.} The extractor was offered every
predicate in the rule base as vocabulary, including rule heads, and duly wrote
\texttt{needs\_supervisor\_review(c)} down as a fact about a case --- a
conclusion asserting itself as evidence. Heads are now withheld.

\textbf{Silent no-ops.} Three times, a string replacement in our own tooling
failed to match and left code that looked correct and did nothing: a prompt
instruction that never existed, a page loader never called, a database column
never written. Each was caught by testing the end state rather than the edit.

\section{Results: coverage and its limiting factor}

\begin{table}[t]
\caption{Two runs of the same twelve-case load from an empty database.}
\label{tab:coverage}
\small
\begin{tabular}{@{}lrr@{}}
\toprule
& \textbf{Run A} & \textbf{Run B} \\
\midrule
Decided from approved logic & 2 / 12 & 3 / 12 \\
Parked for a person         & 5 / 12 & 4 / 12 \\
Coverage                    & 18\%   & 30\%  \\
Approved rules              & 7      & 5     \\
Distinct predicates coined  & 19     & 15    \\
\bottomrule
\end{tabular}
\end{table}

Coverage --- the share of decisions made from approved logic rather than from a
model --- is the measure that matters, and the honest one. It climbs within a
run and does not climb far. The reason is legible in the vocabulary the system
built for itself, which held \texttt{reorder\_toner} and \texttt{reorder\_paper}
side by side: six terms and two rules for one idea.

Coverage appears exactly where the extractor did \emph{not} bake the entity into
the predicate --- \texttt{in\_arrears}, \texttt{debt\_disputed} --- and fails
exactly where it did. Generalisation can widen a rule; it cannot repair a
vocabulary. The engine is first-order and can already express
\texttt{reorder(C, X) $\leftarrow$ stock\_low(C, X)}; nothing proposes it,
because no anti-unification merges two different predicate \emph{names}.

\section{Limitations}

\textbf{No expert has used this.} Every approval in Table~\ref{tab:coverage} was
granted by a script. Whether a reviewer can judge a proposed clause, and how long
it takes them, is the study we have not run and the assumption everything rests
on.

\textbf{Two runs is not a measurement.} The spread between 18 and 30 per cent on
identical input is wide, and driven by whether the extractor happened to reuse a
term.

\textbf{Conduct at the top tier is a prompt.} At \texttt{human\_only} the reply
is constrained by an instruction not to recommend an outcome. It held under
direct pressure in every trial we ran; three trials is an anecdote, and nothing
inspects the reply before it reaches a person.

\textbf{Vocabulary transfers without review.} A rule approved for tenancy fires
on a library fine described in the same terms. That transfer was sound in the
case we observed and is unexamined in general.

\textbf{The matrix has defects we did not fix.} Reversibility does nothing in
the harm 4--7 band, and confidence overrides harm entirely, so a trivially
harmless decision described thinly is sent to a person. Both are properties of
the group's artefact rather than of our implementation, and we report them
rather than silently amending them.

\section{Conclusion}

Human-centred does not mean a person in every loop; at volume that produces
ceremony. It means the unit of human judgement is chosen so that judgement
scales, and that what a person decided is recorded, attributed and reusable.

We have built a system in which the two halves stay apart --- a matrix that
allocates oversight and never sees the answer, a logic database that produces
the answer and never allocates oversight --- joined by a gate whose only
question is whether a person has already ruled on this logic. A fifth of cases
are decided the same way every time, with a derivation attached and a named
person behind the rule. The path from a fifth to a majority runs through a
vocabulary problem we can describe precisely and have not solved.

\end{document}
